\documentclass[11pt]{article}

% Change "review" to "final" to generate the final (sometimes called camera-ready) version.
% Change to "preprint" to generate a non-anonymous version with page numbers.
\usepackage[]{acl}
\usepackage{times}
\usepackage{latexsym}
\usepackage{float} 
\usepackage{placeins}

% For proper rendering and hyphenation of words containing Latin characters (including in bib files)
\usepackage[T1]{fontenc}
% For Vietnamese characters
% \usepackage[T5]{fontenc}
% See https://www.latex-project.org/help/documentation/encguide.pdf for other character sets

% This assumes your files are encoded as UTF8
\usepackage[utf8]{inputenc}
\usepackage{microtype}
\usepackage{inconsolata}
\usepackage{graphicx}

\usepackage{url}
\usepackage{hyperref}

% If the title and author information does not fit in the area allocated, uncomment the following
%
%\setlength\titlebox{<dim>}
%
% and set <dim> to something 5cm or larger.

\title{Group 147 Final Report:\\No More Parties In LA}

\author{Arjun Karthik, Alina Zeng, Stanley Chen \\
  \texttt{\{kartha4,zenga12,chens313\}@mcmaster.ca} }

\begin{document}
\maketitle

\section{Introduction}
% (10%) Introduction
% Copy over and refine the motivation/problem description from progress report.
% Include at least 7 references.
% Discuss related work influencing your direction after the progress report.

\section{Dataset}
% (5%) Dataset
% Describe dataset properties and preprocessing.
% If annotated, describe annotation procedure.
% Describe any changes since progress report (new sources, augmentation, etc.).
% If unchanged, may reuse progress report content.

\section{Features and Inputs}
% (5%) Features and Inputs
% Describe model inputs and feature engineering / representation learning.
% Explain why these features were used.
% Describe any feature selection or augmentation.
% Explain how feature variants were used in experiments.

\section{Implementation}
% (20%) Implementation
% Describe all model versions used.
% Include simple baseline (e.g., majority vote) and trained baselines.
% Explain loss function, optimization technique.
% Include ablations if model has many parts.
% If model underperformed a baseline, explain why.

\section{Evaluation}
% (15%) Evaluation
% Describe evaluation strategy: train/val/test splits, distributions.
% Identify evaluation metrics and justify them.
% Compare to progress report metrics and whether they were adequate.

\section{Progress}
% (10%) Progress
% Reflect on original plan vs. what happened.
% Explain deviations from plan and why.

\section{Error Analysis}
% (20%) Error Analysis
% Systematically analyze model errors.
% Include figures, confusion matrices, qualitative samples, etc.
% Identify patterns in good/bad performance.
% Suggest specific improvements for future work.

\section*{Team Contributions}
% (5%) Team Contributions
% Provide at least one sentence per team member.
% Describe concrete contributions; do not write "all contributed equally".

\begin{itemize}
    \item Alina: Sample text.
    \item Arjun: Sample text.
    \item Stanley: Sample text.
\end{itemize}

\vspace{1em}

\renewcommand{\refname}{}
\vspace{-3em}

\begin{thebibliography}{9}

% Fill in your references here.

\bibitem{FakeBibItem}
My template citation.

\end{thebibliography}

\end{document}